\subsection{Modules over an Artinian Ring}

PROPOSITION 2. --- Let A be a left Artinian ring. For every A-module M, the following properties are equivalent:

\begin{enumerate}
\def\labelenumi{(\roman{enumi})}
\item
  M is semisimple.
\item
  M is without radical.
\item
  M is annihilated by the radical r of A.
\end{enumerate}

We know that (i) implies (ii) (VIII, p.~153, Proposition 3). On the other hand, rM is contained in the radical of M (VIII, p.~158, Proposition 6), so (ii) implies (iii).

Suppose that the A-module M is annihilated by r. We can view it as a module over the ring A/r. Now, the ring A/r is semisimple (VIII, p.~173, Proposition 1), and every module over a semisimple ring is semisimple (VIII, p.~138, Proposition 4). Consequently, M is a semisimple module over the ring A/r and a fortiori over the ring A. Hence (iii) implies (i).

COROLLARY. --- Let A be a left Artinian ring. For every A-module M, the radical of M is equal to rM.

The radical \(\Re(\mathrm{M})\) of M contains rM (VIII, p.~158, Proposition 6). Moreover, the A-module M/rM is annihilated by r. By Proposition 2, it is therefore without radical, which implies \(\Re(\mathrm{M}) \subset \mathfrak{rM}\) (VIII, p.~152, Corollary 1, c)).

PROPOSITION 3. --- Let A and B be rings and f a homomorphism from A to B such that B is generated by the union of \(f(\mathrm{A})\) and the commutant \(f(\mathrm{A})'\) of \(f(\mathrm{A})\) in B. Let M be a left B-module. Suppose that either the ring A is left Artinian or the A-module \(f_{*}(\mathrm{M})\) deduced from M by restriction of scalars (II, §1, No.~13, p.~221) is finitely generated. We then have \(\mathfrak{R}_{\mathrm{A}}(f_{*}(\mathrm{M})) \subset \mathfrak{R}_{\mathrm{B}}(\mathrm{M})\) .

Let S be a simple B-module. Suppose that either A is left Artinian or the A-module \(f_{*}(S)\) is finitely generated. We will prove that \(f_{*}(S)\) is without radical. For every \(b \in f(A)'\) , the homothety \(b_{S}\) is an endomorphism of the A-module \(f_{*}(S)\) and therefore preserves \(\mathfrak{R}_{\mathrm{A}}(f_{*}(\mathrm{S}))\) (VIII, p.~152, Proposition 1). Since B is generated by \(f(\mathrm{A}) \cup f(\mathrm{A}')\) , the radical \(\mathfrak{R}_{\mathrm{A}}(f_{*}(\mathrm{S}))\) is a B-submodule of S, hence equal to 0 or S. If \(f_{*}(S)\) is a finitely generated A-module, then we have \(\mathfrak{R}_{\mathrm{A}}(f_{*}(\mathrm{S})) \neq f_{*}(\mathrm{S})\) (VIII, p.~153, Proposition 2). If the ring A is left Artinian, then we have \(\mathfrak{R}_{\mathrm{A}}(f_{*}(\mathrm{S})) = f(\mathfrak{r})\mathrm{S}\) , where r is the radical of A (VIII, p.~174, Corollary of Proposition 2), and since r is a nilpotent two-sided ideal of A (VIII, p.~173, Proposition 1), we cannot have \(S = f(\mathfrak{r})S\) . We therefore have \(\mathfrak{R}_{\mathrm{A}}(f_{*}(\mathrm{S})) = 0\) in both cases.

Let u be a nonzero B-linear mapping from M to a simple M-module S. The mapping u is surjective; consequently, if \(f_{*}(\mathrm{M})\) is finitely generated, then so is \(f_{*}(\mathrm{S})\) . Under the assumptions of Proposition 3, the A-module \(f_{*}(\mathrm{S})\) is without radical and u is an A-linear mapping from \(f_{*}(\mathrm{M})\) to \(f_{*}(\mathrm{S})\) . Hence, the kernel of u contains \(\mathfrak{R}_{\mathrm{A}}(f_{*}(\mathrm{M}))\) by Proposition 1 of VIII, p.~152. Since u is arbitrary, we have \(\mathfrak{R}_{\mathrm{A}}(f_{*}(\mathrm{M})) \subset \mathfrak{R}_{\mathrm{B}}(\mathrm{M})\) .

COROLLARY. --- Let A be a commutative ring and B an algebra over A. Suppose that either the ring A is Artinian or B is a finitely generated A-module. We then have \(\Re(\mathrm{A})\mathrm{B}\subset\Re(\mathrm{B})\) .

By Proposition 3, the radical \(\mathfrak{R}_{\mathrm{A}}(\mathrm{B}_{s})\) is contained in \(\mathfrak{R}_{\mathrm{B}}(\mathrm{B}_{s})\) , which is nothing but the radical of the ring B. Moreover, we have \(\mathfrak{R}(\mathrm{A})\mathrm{B}\subset\mathfrak{R}_{\mathrm{A}}(\mathrm{B}_{s})\) by Proposition 6 of VIII, p.~158. The corollary follows.

